\section*{Bab \ref{ch:b1:vspaces} --- Ruang Vektor}

\begin{solution}{exo:b1:vspaces:1}
\begin{enumerate}
  \item Ya: karena ia memuat $0$, dan persamaan pendefinisinya linear
        (sehingga stabil terhadap $x + \lambda y$).
  \item Bukan: karena ia tak memuat $(0,0,0)$.
  \item Bukan: karena $(1, 0)$ dan $(0, -1)$ termasuk (dengan $xy = 0$), sedangkan jumlahnya $(1,
        -1)$ tidak (karena $xy = -1 < 0$).
  \item Ya: karena $0$ lenyap di $1$; dan $(P + \lambda Q)(1) = P(1) +
        \lambda Q(1) = 0$.
  \item Ya: karena fungsi nolnya terbatas; dan jika $\abs f \leq M$ dan
        $\abs g \leq M'$, maka $\abs{f + \lambda g} \leq M +
        \abs\lambda M'$.
\end{enumerate}
\end{solution}

\begin{solution}{exo:b1:vspaces:2}
Persamaan $(1,2,1) = a(1,0,1) + b(1,1,0)$ menuntut $a + b = 1$, $b = 2$, $a =
1$: yang tak selaras (karena $a + b = 3 \neq 1$): jadi tak di \omterm{def:b1:vspaces:span}{rentangnya}. Sedangkan $(2,1,1) =
a(1,0,1) + b(1,1,0)$: $b = 1$, $a = 1$, $a + b = 2$: yang taat asas, sehingga
$(2,1,1) = (1,0,1) + (1,1,0)$, jadi di \omterm{def:b1:vspaces:span}{rentangnya}.

Persamaannya: $(x, y, z) = (a + b, b, a)$ berarti $x = y + z$: sehingga \omterm{def:b1:vspaces:span}{rentangnya}
adalah bidang $\{x - y - z = 0\}$.
\end{solution}

\begin{solution}{exo:b1:vspaces:3}
Keluarga pertamanya: $\lambda(1,1,0) + \mu(1,0,1) + \nu(0,1,1) = 0$ memberikan
$\lambda + \mu = 0$, $\lambda + \nu = 0$, $\mu + \nu = 0$: lalu menjumlahkannya,
$2(\lambda + \mu + \nu) = 0$, dan mengurangkan setiap persamaannya,
$\lambda = \mu = \nu = 0$: jadi \omterm{def:b1:vspaces:free}{bebas}.

Yang kedua: $(2,4,6) = 2(1,2,3)$: jadi terikat.

Yang ketiga: empat vektor di $\R^3$ --- yang niscaya terikat begitu dimensinya
tersedia (\cref{ch:b1:findim}); adapun secara langsung: $(0,1,1) = -(1,0,0) +
0\cdot(1,1,0) + (1,1,1)$, memang $(-1,0,0) + (1,1,1) = (0,1,1)$: yaitu sebuah
relasi yang taksepele.
\end{solution}

\begin{solution}{exo:b1:vspaces:4}
\omterm{def:b1:poly:def}{Polinomial} $1, (X-1), (X-1)^2, (X-1)^3$ berderajat berbeda $0,
1, 2, 3$: jadi \omterm{def:b1:vspaces:free}{bebas} (\cref{prop:b1:vspaces:freecriteria}~(1)); dan empat
vektor \omterm{def:b1:vspaces:free}{bebas} yang membangun (karena setiap $P \in \R_3[X]$ terjabarkan dalam pangkat $X -
1$, misalnya lewat Taylor bagi \omterm{def:b1:poly:def}{polinomial}, bandingkan bukti
\cref{prop:b1:poly:multiplicity}): jadi sebuah \omterm{def:b1:vspaces:free}{basis}. Untuk $X^3$, lewat Taylor di
$1$: $P = X^3$, $P(1) = 1$, $P'(1) = 3$, $P''(1) = 6$, $P'''(1) =
6$:
\[
X^3 = 1 + 3(X - 1) + 3(X-1)^2 + (X-1)^3 ,
\]
dengan \omterm{prop:b1:vspaces:coordinates}{koordinatnya} $(1, 3, 3, 1)$ (yaitu baris Pascal, sebagaimana diharapkan dari $X^3 =
((X-1)+1)^3$).
\end{solution}

\begin{solution}{exo:b1:vspaces:5}
Untuk $F \cap G$: syarat $x = y = z$ dan $x = t = 0$ bersama-sama memberikan
$x = 0$, sehingga $y = z = 0$, dan $t = 0$: jadi irisannya adalah
$\{0\}$. Adapun jumlahnya: diberikan $(x,y,z,t)$, carilah $(a,a,a,b) \in F$ dan
$(0,c,d,0) \in G$ yang berjumlah demikian: $a = x$, $b = t$, $c = y - x$, $d =
z - x$: yang selalu mungkin. Jadi $\R^4 = F \oplus G$, dan
\[
(1,2,3,4) = (1,1,1,4) + (0,1,2,0) .
\]
\end{solution}

\begin{solution}{exo:b1:vspaces:6}
Sebuah unsur $G$ berbentuk $\lambda u$; dan jika ia konvergen, maka (karena
$\lambda u_n = \lambda(-1)^n$ mempunyai dua limit \omterm{def:b1:seq:subsequence}{subbarisan} $\pm
\lambda$) niscaya $\lambda = 0$: sehingga $F \cap G = \{0\}$.

Adapun $F + G$ bukan segalanya: karena ia terdiri atas barisan berbentuk
$c_n + \lambda(-1)^n$ dengan $(c_n)$ yang konvergen. Sedangkan barisan $v_n =
n$ tak berbentuk demikian (karena $v_n - \lambda(-1)^n$ tak terbatas, sehingga tak pernah
konvergen). Jadi $F \oplus G \subsetneq$ (ruang semua barisannya).
\end{solution}

\begin{solution}{exo:b1:vspaces:7}
($\supseteq$) Baik $F \cap G$ maupun $F \cap H$ terletak di $F$, dan
jumlahnya terletak di $G + (F \cap H)$: sehingga inklusinya menyusul karena ruas kirinya
merupakan \omterm{def:b1:vspaces:subspace}{subruang} yang memuat kedua kepingnya --- secara konkret, sebuah unsur $g +
h$ dengan $g \in F\cap G$, $h \in F \cap H$ berada di $F$ (sebagai jumlah dua
unsur $F$) dan di $G + (F \cap H)$.

($\subseteq$) Misalkan $x \in F$ dengan $x = g + h$, $g \in G$, $h \in F
\cap H$. Maka $g = x - h \in F$ (sebagai selisih unsur $F$), sehingga
$g \in F \cap G$, dan $x = g + h \in (F \cap G) + (F \cap H)$.

Adapun contoh penyangkal bagi kedistributifan penuhnya di $\R^2$: $F =
\operatorname{Vect}(1,1)$, $G = \operatorname{Vect}(1,0)$, $H =
\operatorname{Vect}(0,1)$. Maka $G + H = \R^2$, sehingga $F \cap (G+H) =
F$, sedangkan $F \cap G = F \cap H = \{0\}$: jadi ruas kanannya $\{0\}
\neq F$.
\end{solution}

\begin{solution}{exo:b1:vspaces:8}
\begin{enumerate}
  \item Andaikan $\sum_{i} \lambda_i \eu^{a_i x} = 0$ untuk setiap $x$.
        Kalikanlah dengan $\eu^{-a_p x}$: maka $\lambda_p + \sum_{i < p}
        \lambda_i \eu^{(a_i - a_p)x} \to \lambda_p$ ketika $x \to
        +\infty$ (karena setiap eksponennya $a_i - a_p < 0$). Padahal ruas kirinya
        identik $0$, sehingga $\lambda_p = 0$; lalu ulangilah ke bawah.
  \item Misalkan $a\cos x + b \sin x + c \cos 2x + d\sin 2x = 0$ untuk setiap
        $x$. Nilailah di $x = 0$: $a + c = 0$; lalu di $x = \pi$: $-a + c
        = 0$; sehingga $a = c = 0$, dan relasinya tereduksi menjadi $b\sin
        x + d \sin 2x = 0$. Lalu nilailah di $x = \frac\pi2$: $b = 0$;
        kemudian di $x = \frac\pi4$: $d = 0$.
  \item Andaikan $\sum \lambda_i \abs{x - a_i} = 0$ untuk setiap $x$. Adapun
        fungsi $\sum_{i \neq j} \lambda_i\abs{x - a_i}$ bersifat
        \omterm{def:b1:derivative:def}{dapat diturunkan} di $a_j$ (karena setiap sukunya demikian, jauh dari
        sudutnya sendiri), sehingga $-\lambda_j \abs{x - a_j}$, yaitu selisihnya,
        pun harus \omterm{def:b1:derivative:def}{dapat diturunkan} di $a_j$ --- yang memaksa
        $\lambda_j = 0$ (karena $\abs{\,\cdot\,}$ bersudut). Dan ini berlaku
        bagi setiap $j$.
\end{enumerate}
\end{solution}

\begin{solution}{exo:b1:vspaces:9}
Misalkan $h \in H$. Karena $h \in H \subseteq F + H = F + G$, tulislah $h = f
+ g$ dengan $f \in F$, $g \in G$. Maka $f = h - g \in H$ (karena kedua sukunya
di $H$, dengan memakai $G \subseteq H$), sehingga $f \in F \cap H = F \cap G
\subseteq G$, dan $h = f + g \in G$. Jadi $H \subseteq G$, lalu digabung dengan
hipotesis $G \subseteq H$: sama.

Adapun contoh penyangkal tanpa $G \subseteq H$: di $\R^2$, ambillah $F =
\operatorname{Vect}(1,0)$, $G = \operatorname{Vect}(0,1)$, $H =
\operatorname{Vect}(1,1)$: maka $F + G = F + H = \R^2$ dan $F \cap G
= F \cap H = \{0\}$, namun $G \neq H$.
\end{solution}

\begin{solution}{exo:b1:vspaces:10}
Kedua \omterm{def:b1:logic:sets}{himpunannya} merupakan \omterm{def:b1:vspaces:subspace}{subruang} (karena syarat pendefinisinya linear dan
berlaku bagi $0$). Adapun penguraiannya: untuk $P \in \R[X]$,
\[
P(X) = \underbrace{\frac{P(X) + P(-X)}{2}}_{\in\,\mathcal P}
+ \underbrace{\frac{P(X) - P(-X)}{2}}_{\in\,\mathcal I},
\]
dan sebuah \omterm{def:b1:poly:def}{polinomial} yang sekaligus genap dan ganjil memenuhi $P = -P$, sehingga $P = 0$:
jadi jumlahnya langsung dan sama dengan $\R[X]$.

Sekarang misalkan $P = \sum_k a_k X^k$ genap. Maka $P(X) - P(-X) = 2
\sum_{k \text{ ganjil}} a_k X^k$ merupakan \omterm{def:b1:poly:def}{polinomial} nol, sehingga setiap
koefisien berderajat ganjilnya lenyap (\cref{def:b1:poly:def}): jadi $P \in
\operatorname{Vect}(1, X^2, X^4, \dots)$, yakni $P = Q(X^2)$ untuk suatu
\omterm{def:b1:poly:def}{polinomial} $Q$. Sebaliknya setiap \omterm{def:b1:poly:def}{polinomial} dalam $X^2$ bersifat genap.
\end{solution}

\begin{solution}{exo:b1:vspaces:11}
Andaikan $a(x_1 + x_2) + b(x_2 + x_3) + c(x_3 + x_1) = 0$.
Lalu mengelompokkannya ulang pada \omterm{def:b1:vspaces:free}{keluarga bebas} $(x_1, x_2, x_3)$:
\[
(a + c)\,x_1 + (a + b)\,x_2 + (b + c)\,x_3 = 0
\implies a + c = a + b = b + c = 0 .
\]
Mengurangkan kedua persamaan pertamanya memberikan $c = b$; lalu yang ketiga
memberikan $2b = 0$, sehingga $b = c = 0$, kemudian $a = 0$: jadi keluarganya \omterm{def:b1:vspaces:free}{bebas}.

Adapun untuk empat vektor keluarga serupanya \emph{selalu} terikat:
\[
(x_1 + x_2) - (x_2 + x_3) + (x_3 + x_4) - (x_4 + x_1) = 0
\]
merupakan kombinasi nol yang taksepele (dengan koefisien $1, -1, 1, -1$),
apa pun $(x_1, x_2, x_3, x_4)$-nya. Jadi keparitasan panjang daurnyalah
yang memutuskan.
\end{solution}

\begin{solution}{exo:b1:vspaces:12}
\begin{enumerate}
  \item Jika $F_1 \subseteq F_2$ atau $F_2 \subseteq F_1$, maka gabungannya
        salah satu dari keduanya, sehingga sejati. Kalau tidak pilihlah $x \in F_1
        \setminus F_2$ dan $y \in F_2 \setminus F_1$, lalu tinjaulah
        $x + y$. Jika $x + y \in F_1$, maka $y = (x + y) - x \in
        F_1$: yang bertentangan. Jika $x + y \in F_2$, maka $x \in F_2$:
        yang bertentangan. Jadi $x + y \notin F_1 \cup F_2$, sehingga $E \neq
        F_1 \cup F_2$.
  \item Andaikan untuk pertentangannya bahwa $E = F_1 \cup \dots \cup
        F_k$, dengan $k$ yang dipilih \emph{minimal} di antara semua
        penyelimutan yang demikian. Maka keminimalannya melarang $F_1 \subseteq F_2 \cup
        \dots \cup F_k$ (kalau tidak buanglah $F_1$), sehingga ada $x \in F_1$
        dengan $x \notin F_i$ untuk setiap $i \geq 2$. Dan karena $F_1$
        sejati, pilihlah $y \notin F_1$. Untuk setiap skalar $t$,
        vektor $y + t x$ terletak di suatu $F_i$. Namun ia tak pernah terletak di
        $F_1$: sebab kalau tidak $y = (y + tx) - tx \in F_1$ (karena $x \in
        F_1$). Adapun \omterm{def:b1:structures:field}{lapangannya} tak hingga, sehingga pilihlah $k$ skalar
        yang berbeda $t_1, \dots, t_k$: maka $k$ vektor $y + t_j x$
        jatuh ke dalam $k - 1$ \omterm{def:b1:vspaces:subspace}{subruang} $F_2, \dots, F_k$, sehingga dua
        di antaranya, katakanlah $y + t x$ dan $y + t' x$ dengan $t \neq t'$,
        terletak di $F_i$ yang sama (dengan $i \geq 2$). Lalu selisihnya
        $(t - t')x \in F_i$, sehingga $x \in F_i$: yang bertentangan. Jadi
        tak ada penyelimutan hingga oleh \omterm{def:b1:vspaces:subspace}{subruang} sejati.
\end{enumerate}
\end{solution}

\begin{solution}{pb:b1:vspaces:1}
\textbf{1.} Di sini $L_i$ merupakan hasil kali $n$ faktor linear dibagi
konstanta taknol (karena $x_i$-nya berbeda), sehingga $\deg L_i = n$.
Lalu menilainya di $x_j$ dengan $j \neq i$: faktor $X - x_j$ pada
pembilangnya lenyap, sehingga $L_i(x_j) = 0$. Sedangkan di $x_i$, pembilang dan
penyebutnya berimpit: sehingga $L_i(x_i) = 1$.

\textbf{2.} Andaikan $\sum_i \lambda_i L_i = 0$. Nilailah di $x_j$:
maka semua sukunya mati kecuali $\lambda_j L_j(x_j) = \lambda_j$, sehingga
$\lambda_j = 0$ untuk setiap $j$: jadi keluarganya \omterm{def:b1:vspaces:free}{bebas}.

\textbf{3.} Misalkan $D = P - \sum_i P(x_i) L_i$. Maka $\deg D \leq n$
dan, menurut pertanyaan 1, $D(x_j) = P(x_j) - P(x_j) = 0$ bagi $n + 1$
titik berbeda $x_0, \dots, x_n$. Adapun \omterm{def:b1:poly:def}{polinomial} taknol berderajat
$\leq n$ berakar paling banyak $n$ (\cref{cor:b1:poly:nroots}), sehingga $D
= 0$. Jadi setiap $P \in \R_n[X]$ merupakan kombinasi $L_i$-nya:
sehingga keluarganya pembangun, dan digabung dengan pertanyaan 2, sebuah \omterm{def:b1:vspaces:free}{basis}.

\textbf{4.} Diberikan $y_0, \dots, y_n$, \omterm{def:b1:poly:def}{polinomial} $P = \sum_i
y_i L_i$ berderajat $\leq n$ dan menginterpolasinya. Adapun ketunggalannya: sebuah
$P$ yang menginterpolasi mempunyai, menurut pertanyaan 3, \omterm{prop:b1:vspaces:coordinates}{koordinat} $(P(x_0), \dots,
P(x_n)) = (y_0, \dots, y_n)$ pada \omterm{def:b1:vspaces:free}{basis} $(L_i)$, dan \omterm{prop:b1:vspaces:coordinates}{koordinat}
pada sebuah \omterm{def:b1:vspaces:free}{basis} bersifat tunggal (\cref{prop:b1:vspaces:coordinates}). Jadi
\omterm{prop:b1:vspaces:coordinates}{koordinat} $P$ pada \omterm{def:b1:vspaces:free}{basis} Lagrangenya adalah \emph{nilainya di
simpulnya} --- dan itulah seluruh inti basisnya.

\textbf{5.} Terapkanlah pertanyaan 3 pada $P = X^k$ (dengan $0 \leq k \leq n$):
\[
X^k = \sum_{i=0}^{n} x_i^{k} L_i ,
\]
lalu $k = 0$ memberikan $\sum_i L_i = 1$.

\textbf{6.} Derajat $\deg N_k = k$ persis: sehingga keluarga $(N_0, \dots, N_n)$
merupakan tangga derajat di $\R_n[X]$, jadi sebuah \omterm{def:b1:vspaces:free}{basis} menurut
\cref{ex:b1:vspaces:staircase} (dengan kebebasannya dari
\cref{prop:b1:vspaces:freecriteria}~(1), dan pembangunnya lewat penurunan
hingga pada derajatnya).

\textbf{7.} Dari rekurensinya,
\[
f[x_0, x_1] = \frac{f(x_1) - f(x_0)}{x_1 - x_0},
\qquad
f[x_0, x_1, x_2] = \frac{f[x_1, x_2] - f[x_0, x_1]}{x_2 - x_0}.
\]
Untuk $f(x) = x^2$:
\[
f[x_0, x_1] = \frac{x_1^2 - x_0^2}{x_1 - x_0} = x_0 + x_1,
\]
lalu kemudian
\[
f[x_0, x_1, x_2]
= \frac{(x_1 + x_2) - (x_0 + x_1)}{x_2 - x_0}
= \frac{x_2 - x_0}{x_2 - x_0} = 1 .
\]

\textbf{8.} Derajat $\deg S \leq n$ karena $Q, R$ berderajat $\leq n - 1$.
Di $x_0$: $S(x_0) = \frac{-(x_0 - x_n) R(x_0)}{x_n - x_0} =
R(x_0) = f(x_0)$. Di $x_n$: $S(x_n) = \frac{(x_n - x_0)
Q(x_n)}{x_n - x_0} = Q(x_n) = f(x_n)$. Sedangkan di sebuah simpul \omterm{def:b1:topology:closure}{interior} $x_i$
(dengan $1 \leq i \leq n-1$), baik $Q$ maupun $R$ bernilai $f(x_i)$,
sehingga
\[
S(x_i) = \frac{(x_i - x_0) - (x_i - x_n)}{x_n - x_0}\, f(x_i)
= f(x_i) .
\]

\textbf{9.} Lewat induksi pada cacah simpulnya. Untuk satu simpul:
interpolannya adalah konstanta $f(x_0) = f[x_0]$. Anggaplah klaimnya berlaku
bagi $k$ simpul lalu misalkan $S$ menginterpolasi di $x_0, \dots, x_k$; maka menurut
ketunggalannya (pertanyaan 4), $S$ diberikan oleh lema Aitken dari $R$
(dengan simpul $x_0, \dots, x_{k-1}$) dan $Q$ (dengan simpul $x_1, \dots, x_k$).
Adapun koefisien $X^{k}$ pada $S$ adalah
\[
\frac{[X^{k-1}]\,Q - [X^{k-1}]\,R}{x_k - x_0}
= \frac{f[x_1, \dots, x_k] - f[x_0, \dots, x_{k-1}]}{x_k - x_0}
= f[x_0, \dots, x_k]
\]
menurut hipotesis induksinya dan rekurensi pendefinisinya.

\textbf{10.} Misalkan $P_k$ menginterpolasi $f$ di $x_0, \dots, x_k$. Maka
selisihnya $P_k - P_{k-1}$ berderajat $\leq k$ dan lenyap di
$x_0, \dots, x_{k-1}$, sehingga menurut teorema faktor yang diterapkan $k$ kali
(\cref{thm:b1:poly:factor}) ia sama dengan $c\,N_k$ untuk suatu konstanta $c$;
lalu membandingkan koefisien $X^{k}$-nya dan memakai pertanyaan 9, $c =
f[x_0, \dots, x_k]$. Lalu berteleskop dari $P_0 = f(x_0) N_0$ memberikan
rumus Newtonnya. Adapun untuk bentuk \omterm{def:b1:topology:closed}{tertutupnya}, tulislah $P_k = \sum_{i \leq
k} f(x_i) L_i$ (yaitu Lagrange, pada simpul $x_0, \dots, x_k$) lalu bacalah
koefisien $X^{k}$-nya: karena setiap $L_i$ menyumbang
$\frac{1}{\prod_{j \neq i}(x_i - x_j)}$, sehingga
\[
f[x_0, \dots, x_k] = \sum_{i=0}^{k}
\frac{f(x_i)}{\prod_{j \neq i,\, j \leq k}(x_i - x_j)} .
\]
Adapun ruas kanannya kekal di bawah sebarang \omterm{def:b1:counting:objects}{permutasi} simpulnya, sehingga
\omterm{pb:b1:vspaces:1}{selisih terbaginya} tak bergantung pada urutannya.

\textbf{11.} Jika $P = a X^m + (\text{derajat yang lebih rendah})$, maka teorema
binomialnya memberikan
\[
\Delta P = a\bigl((X+1)^m - X^m\bigr) + \dots
= a\,m\,X^{m-1} + (\text{derajat yang lebih rendah}),
\]
karena $(X+1)^m - X^m = m X^{m-1} + \dots$ dan bagian berderajat rendah
$P$ menyumbang derajat $\leq m - 2$ setelah $\Delta$ (atau
suku berderajat $\leq m-2$). Jadi $\deg \Delta P = m - 1$ dengan koefisien
utamanya $m a$. Adapun sebuah konstanta $c$ memberikan $\Delta c = c - c = 0$.

\textbf{12.} Derajat $\deg B_k = k$: jadi sebuah tangga, sehingga sebuah \omterm{def:b1:vspaces:free}{basis}
$\R_n[X]$ (\cref{ex:b1:vspaces:staircase}). Adapun untuk $\Delta B_k$
(dengan $k \geq 1$), faktorkanlah hasil kali bersamanya:
\begin{align*}
k!\,\Delta B_k
&= (X+1)X\cdots(X-k+2) - X(X-1)\cdots(X-k+1) \\
&= X(X-1)\cdots(X-k+2)\,\bigl[(X+1) - (X-k+1)\bigr] \\
&= k\,X(X-1)\cdots(X-k+2),
\end{align*}
sehingga $\Delta B_k = \frac{X(X-1)\cdots(X-k+2)}{(k-1)!} = B_{k-1}$.

\textbf{13.} Tulislah $P = \sum_{k=0}^{n} c_k B_k$ (menurut basisnya, pertanyaan
12). Lalu terapkanlah $\Delta^{j}$: menurut pertanyaan 12, $\Delta^{j} P = \sum_{k
\geq j} c_k B_{k-j}$. Lalu nilailah di $0$: $B_0(0) = 1$ dan $B_m(0) =
0$ untuk $m \geq 1$ (karena faktor $X$-nya lenyap), sehingga
$\bigl(\Delta^{j}P\bigr)(0) = c_j$. Inilah rumus selisih
majunya.

\textbf{14.} Lewat induksi pada $k$. Untuk $k = 0$ kesamaannya berbunyi
$P(0) = P(0)$. Anggaplah ia berlaku bagi $k$ lalu terapkanlah pada $\Delta P$:
\[
\bigl(\Delta^{k+1} P\bigr)(0)
= \sum_{j=0}^{k} (-1)^{k-j}\binom kj \bigl(P(j+1) - P(j)\bigr).
\]
Lalu kumpulkanlah koefisien $P(i)$-nya: yaitu $(-1)^{k-i+1}\binom
k{i-1}\cdot(-1)^{0}$ dari jumlah pertamanya (yang tergeser) dan
$-(-1)^{k-i}\binom ki$ dari yang kedua --- sehingga bersama-sama
\[
(-1)^{k+1-i}\Bigl(\binom k{i-1} + \binom ki\Bigr)
= (-1)^{k+1-i}\binom{k+1}i
\]
menurut aturan Pascal, yang merupakan kesamaannya pada peringkat $k + 1$.

\textbf{15.} Dengan mengulangi pertanyaan 11 dari derajat $n$, dengan koefisien
utamanya $a_n$: setelah satu $\Delta$, derajatnya $n-1$ dan koefisien
utamanya $n a_n$; setelah dua, $n(n-1)a_n$; setelah $n$ langkah,
derajatnya $0$ dan nilainya $n(n-1)\cdots 1\, a_n = n!\,a_n$, yaitu konstanta.
Lalu satu $\Delta$ lagi membunuhnya: $\Delta^{n+1}P = 0$.

\textbf{16.} Jika $m \geq k$: maka $B_k(m) = \binom mk \in \N$. Jika $0
\leq m < k$: maka satu faktor $m(m-1)\cdots(m-k+1)$ bernilai nol, sehingga
$B_k(m) = 0$. Sedangkan jika $m = -q$ dengan $q \geq 1$:
\[
B_k(-q) = \frac{(-q)(-q-1)\cdots(-q-k+1)}{k!}
= (-1)^k\,\frac{q(q+1)\cdots(q+k-1)}{k!}
= (-1)^k \binom{q+k-1}{k},
\]
yang bulat. Jadi setiap $B_k$ memetakan $\Z$ ke dalam $\Z$.

\textbf{17.} ($\Leftarrow$) Jika $P = \sum_k c_k B_k$ dengan $c_k \in
\Z$, maka untuk $m \in \Z$, $P(m) = \sum_k c_k B_k(m) \in \Z$ menurut
pertanyaan 16. ($\Rightarrow$) Jika $P$ bernilai bulat, maka
\omterm{prop:b1:vspaces:coordinates}{koordinatnya} adalah $c_k = \bigl(\Delta^k P\bigr)(0) = \sum_{j=0}^k
(-1)^{k-j}\binom kj P(j)$ (menurut pertanyaan 13 dan 14), yaitu kombinasi
bulat atas bilangan bulat $P(0), \dots, P(k)$. Inilah
pencirian Pólya atas \omterm{pb:b1:vspaces:1}{polinomial bernilai bulat}.

\textbf{18.} Ambillah $Q(X) = P(X + a)$, \omterm{def:b1:poly:def}{polinomial}
berderajat $\leq n$, dengan $Q(0), Q(1), \dots, Q(n) \in \Z$. Adapun \omterm{prop:b1:vspaces:coordinates}{koordinatnya} pada
$(B_k)_{k \leq n}$ adalah $c_k = \sum_{j \leq k}(-1)^{k-j}\binom kj
Q(j) \in \Z$ (karena pertanyaan 14 hanya memakai nilainya di $0, \dots, k
\leq n$). Jadi menurut pertanyaan 17 ($\Leftarrow$), $Q$ bernilai bulat pada
seluruh $\Z$, sehingga $P(X) = Q(X - a)$ pun demikian.

\textbf{19.} Hasil kali $k$ bilangan bulat yang berurutan adalah $m(m-1)
\cdots(m-k+1) = k!\,B_k(m)$ untuk suatu $m \in \Z$, dan $B_k(m) \in
\Z$ menurut pertanyaan 16: sehingga hasil kalinya terbagi oleh $k!$.

\textbf{20.} Nilai $P = \frac{X(X+1)(2X+1)}{6}$ di $0,1,2,3$:
$0, 1, 5, 14$. Adapun tabel selisihnya: baris $\Delta$-nya $1, 4, 9$;
baris $\Delta^2$-nya $3, 5$; dan baris $\Delta^3$-nya $2$. Sehingga, menurut pertanyaan 13,
\[
P = 0\cdot B_0 + 1\cdot B_1 + 3\,B_2 + 2\,B_3 ,
\]
dengan \omterm{prop:b1:vspaces:coordinates}{koordinat} yang bulat: sehingga $P$ bernilai bulat (pertanyaan 17),
sedangkan koefisien monomialnya $\frac13, \frac12, \frac16$
tidak bulat. Adapun perhitungan langsungnya:
\begin{align*}
\Delta P &= \frac{(X+1)(X+2)(2X+3) - X(X+1)(2X+1)}{6} \\
&= \frac{(X+1)\bigl[(X+2)(2X+3) - X(2X+1)\bigr]}{6}
= \frac{(X+1)(6X+6)}{6} = (X+1)^2 .
\end{align*}
Lalu berteleskop $P(m) = \sum_{j=0}^{m-1}\Delta P(j) = \sum_{j=1}^{m}
j^2$ (dengan $P(0) = 0$): yaitu rumus jumlah kuadratnya.

\textbf{21.} Nilai $2^i$ di $i = 0, \dots, n$ mempunyai tabel
selisih yang selalu $1$ pada tepi kirinya: karena $\Delta^k$ atas barisan
$(2^i)$ kembali menjadi $(2^i)$ (sebab $2^{i+1} - 2^i = 2^i$), sehingga
$\bigl(\Delta^k P\bigr)(0) = 2^0 = 1$ untuk setiap $k \leq n$, dan $P =
B_0 + B_1 + \dots + B_n$ menurut pertanyaan 13. Lalu
\[
P(n+1) = \sum_{k=0}^{n}\binom{n+1}{k}
= 2^{n+1} - \binom{n+1}{n+1} = 2^{n+1} - 1 \neq 2^{n+1}:
\]
sehingga polanya patah pada titik pertama yang tak terkendali.

\textbf{22.} Menurut pertanyaan 12, $B_k = \Delta B_{k+1}$, sehingga
\[
\sum_{j=0}^{m-1} B_k(j)
= \sum_{j=0}^{m-1}\bigl(B_{k+1}(j+1) - B_{k+1}(j)\bigr)
= B_{k+1}(m) - B_{k+1}(0) = B_{k+1}(m).
\]
Adapun untuk $j < k$ suku $B_k(j)$-nya lenyap, sehingga jumlahnya sungguh bermula di
$j = k$: $\sum_{j=k}^{m-1}\binom jk = \binom m{k+1}$, yaitu
kesamaan tongkat hokinya.

\textbf{23.} Tabel selisihnya (atau penjabaran langsungnya) memberikan
\[
X^2 = B_1 + 2 B_2, \qquad X^3 = B_1 + 6 B_2 + 6 B_3
\]
(periksa: $B_1 + 2B_2 = X + X(X-1) = X^2$; dan di $X = 1, 2, 3$ yang
kedua memberikan $1, 8, 27$). Lalu pertanyaan 22 menghasilkan
\[
\sum_{j=0}^{m-1} j^2 = B_2(m) + 2B_3(m)
= \binom m2 + 2\binom m3 = \frac{m(m-1)(2m-1)}{6},
\]
\[
\sum_{j=0}^{m-1} j^3 = B_2(m) + 6B_3(m) + 6B_4(m)
= \binom m2 + 6\binom m3 + 6\binom m4 .
\]
Lalu menjabarkan ungkapan terakhirnya: $\binom m2 + 6\binom m3 + 6\binom
m4 = \frac{m(m-1)}{2}\bigl[1 + 2(m-2) +
\frac{(m-2)(m-3)}{2}\bigr] = \frac{m^2(m-1)^2}{4} = \binom m2^2$.
Lalu mengganti $m$ dengan $m + 1$: $1^3 + \dots + m^3 =
\bigl(\frac{m(m+1)}2\bigr)^2 = (1 + \dots + m)^2$, yaitu kesamaan
Nicomachus.

\textbf{24.} Dengan simpul $0, 1, 2$, dan \omterm{def:b1:poly:def}{polinomial} $X^2$. Pada \omterm{def:b1:vspaces:free}{basis} monomialnya
$(1, X, X^2)$: \omterm{prop:b1:vspaces:coordinates}{koordinatnya} $(0, 0, 1)$. Pada \omterm{def:b1:vspaces:free}{basis} Lagrangenya:
\omterm{prop:b1:vspaces:coordinates}{koordinatnya} adalah nilainya $(0, 1, 4)$ (pertanyaan 4). Pada \omterm{def:b1:vspaces:free}{basis} Newtonnya
$(1, X, X(X-1))$: \omterm{pb:b1:vspaces:1}{selisih terbaginya} $f[0] = 0$, $f[0,1] = 1$,
$f[0,1,2] = \frac{3 - 1}{2} = 1$ (pertanyaan 9), sehingga \omterm{prop:b1:vspaces:coordinates}{koordinatnya}
$(0, 1, 1)$ --- memang $X + X(X-1) = X^2$. Jadi tiga \omterm{def:b1:vspaces:free}{basis}, tiga
vektor \omterm{prop:b1:vspaces:coordinates}{koordinat}, satu \omterm{def:b1:poly:def}{polinomial}.

\textbf{25.} (i) Pertanyaan 4 bersifat otomatis karena $(L_i)$ merupakan
sebuah \emph{\omterm{def:b1:vspaces:free}{basis}}: sehingga keberadaan dan ketunggalan interpolasinya
persis merupakan keberadaan dan ketunggalan \omterm{prop:b1:vspaces:coordinates}{koordinatnya}. (ii) Adapun \omterm{def:b1:vspaces:free}{basis}
Newtonnya sebuah tangga, sehingga \omterm{prop:b1:vspaces:coordinates}{koordinatnya} dihitung lewat pembagian yang
berturut-turut --- karena setiap simpul barunya menambahkan satu suku tanpa mengganggu
yang sebelumnya --- sedangkan \omterm{prop:b1:vspaces:coordinates}{koordinat} Lagrange $P$ adalah
nilai $P(x_i)$, yang tersedia tanpa perhitungan sama sekali. (iii) Kedua
basisnya \omterm{def:b1:vspaces:free}{bebas} menurut kedua kriteria yang sama pada
\cref{prop:b1:vspaces:freecriteria}: yaitu derajat yang berbeda bagi Newton,
dan penilaian di simpulnya bagi Lagrange. (iv) Adapun teorema Pólya mengatakan
bahwa ``$P(\Z) \subseteq \Z$'', yaitu sifat nilainya, setara
dengan kebulatan \omterm{prop:b1:vspaces:coordinates}{koordinatnya} pada \omterm{def:b1:vspaces:free}{basis} $(B_k)$ --- sehingga
aritmetika sebuah \omterm{def:b1:poly:def}{polinomial} menjadi kasatmata hanya pada \omterm{def:b1:vspaces:free}{basis}
yang disesuaikan dengan pertanyaannya.

\end{solution}
